% ============================================================
% BÖLÜM 4: ARAŞTIRMANIN ÖNEMİ
% ============================================================

\chapter{ARAŞTIRMANIN ÖNEMİ}

\section{PDR Alanına Katkısı}

Bu çalışma, psikolojik danışmanlık ve rehberlik alanına çeşitli boyutlarda katkı sunmaktadır:

\subsection{Kuramsal Katkılar}

\begin{enumerate}
  \item \textbf{Teknoloji-PDR Entegrasyonu:} Çalışma, yapay zekâ teknolojilerinin psikolojik danışmanlık hizmetlerine nasıl entegre edilebileceğine ilişkin bir model sunmaktadır. Bu model, gelecekteki benzer çalışmalar için referans niteliği taşımaktadır. Luxton (2014), yapay zekânın ruh sağlığı hizmetlerinde dönüştürücü bir potansiyele sahip olduğunu; ancak bu entegrasyonun dikkatli biçimde planlanması gerektiğini vurgulamıştır.

  \item \textbf{Projektif Testlerin Dijitalleşmesi:} HTP gibi geleneksel değerlendirme araçlarının dijital platformlara taşınması, alana yeni bir perspektif kazandırmaktadır. Bordin (1979), terapötik ittifak kavramını tanımlarken teknolojinin bu ittifakı destekleyebileceğini öngörmemiş olsa da, günümüzde dijital araçlar bu ittifakı güçlendirebilmektedir (Berger, 2017).

  \item \textbf{Bilgi Tabanı Sistematizasyonu:} Dağınık akademik literatürün yapılandırılmış bir bilgi tabanına dönüştürülmesi, alana metodolojik bir katkı sunmaktadır. Bu yaklaşım, ``kanıta dayalı uygulama'' (evidence-based practice) paradigması ile uyumludur (Sackett ve ark., 1996).
\end{enumerate}

\subsection{Uygulamaya Yönelik Katkılar}

\begin{enumerate}
  \item \textbf{Verimlilik Artışı:} Okul psikolojik danışmanlarının değerlendirme süreçlerindeki iş yükünün azaltılması hedeflenmektedir. Automated clinical decision support systems have been shown to reduce clinician workload by 20-30\% in healthcare settings (Bright ve ark., 2012).

  \item \textbf{Standartlaşma:} Yorumlamaların akademik literatüre dayandırılması, değerlendirmeler arası tutarlılığı artırma potansiyeli taşımaktadır. Handler (1996), projektif testlerde standartlaşmanın güvenilirliği artırdığını belirtmiştir.

  \item \textbf{Erişilebilirlik:} Web tabanlı yapı, coğrafi konumdan bağımsız erişim imkânı sunmaktadır. Bu özellik, özellikle kırsal bölgelerdeki okullarda görev yapan danışmanlar için önemlidir.

  \item \textbf{Sürekli Güncellenebilirlik:} Dijital platform, yeni akademik bulguların sisteme hızlıca entegre edilmesini mümkün kılmaktadır.
\end{enumerate}

\section{Psikolojik Danışmanlara ve Rehber Öğretmenlere Katkısı}

\begin{itemize}
  \item Çizim yorumlama konusunda anlık destek ve referans kaynağı
  \item Akademik literatüre kolay erişim ve kaynak takibi
  \item Mesleki gelişimi destekleyen eğitici içerikler
  \item Zaman tasarrufu sağlayan otomatik analiz
  \item Raporlama süreçlerinin kolaylaştırılması
  \item Vaka tartışmalarında kullanılabilecek dokümantasyon
  \item Süpervizyon süreçlerinde yardımcı araç
\end{itemize}

\section{Eğitim Ortamlarına Olası Etkileri}

\begin{enumerate}
  \item \textbf{Erken Tespit:} Psikolojik sorunların daha erken dönemde fark edilmesi, zamanında müdahale imkânı sunmaktadır. Werner ve Smith (2001), erken müdahalenin çocukların psikolojik dayanıklılığını önemli ölçüde artırdığını göstermiştir.

  \item \textbf{Kaynakların Etkin Kullanımı:} Sınırlı psikolojik danışman kaynağının daha verimli kullanılmasına katkı sağlama potansiyeli bulunmaktadır.

  \item \textbf{Veri Temelli Karar Alma:} Sistematik kayıt ve raporlama, okul yönetimlerinin veri temelli kararlar almasını destekleyebilir. Bu yaklaşım, ``data-driven decision making'' (DDDM) paradigması ile uyumludur (Mandinach, 2012).

  \item \textbf{Farkındalık Artışı:} Sistemin kullanımı, çocuk ruh sağlığı konusunda eğitim camiasında farkındalık yaratabilir.
\end{enumerate}

\section{Toplumsal Etki}

\begin{enumerate}
  \item \textbf{Ruh Sağlığı Hizmetlerine Erişim:} Türkiye'de ruh sağlığı uzmanı sayısı sınırlıdır. OECD (2021) verilerine göre, Türkiye'de 100.000 kişiye düşen psikiyatrist sayısı OECD ortalamasının altındadır. NIDA gibi destek sistemleri, bu açığı kısmen kapatabilir.

  \item \textbf{Damgalamanın Azaltılması:} Dijital araçlar, ruh sağlığı değerlendirmelerinin ``normalleştirilmesine'' katkı sağlayabilir (Corrigan ve ark., 2014).

  \item \textbf{Önleyici Ruh Sağlığı:} Erken tespit, tedavi maliyetlerini düşürmekte ve bireylerin yaşam kalitesini artırmaktadır.
\end{enumerate}

\newpage
